\clearpage
\fancyhead[L]{\small KOMET.}
\hypertarget{native-474-lower}{}
\begin{center}
\rule{0.55\linewidth}{0.4pt}\\[2pt]
Astronomische Nachrichten. Band XXII. Nr. 521. S. 277. 278. 28. November 1844.\\[2pt]
\rule{0.55\linewidth}{0.4pt}
\end{center}
\begin{flushright}Göttingen, 15. Oktober und 11. November 1844.\end{flushright}

Am 6. Oktober habe ich den \emph{Kometen} ein paarmal mit Ceti 200 \textsc{Bode} verglichen. Aus diesen Vergleichen lässt sich nur die Rektaszension mit Sicherheit ableiten.

\begin{center}
1844 Oktober 6.\quad $10^{\rm u}\ 40^{\rm m}\ 32^{\rm s}$ m. Z.\quad Rektasz. $18^\circ\ 50'\ 51''1$
\end{center}

Für die Sterne, mit denen ich den Kometen am 14. Oktober verglichen habe, lege ich aufgrund der verschiedenen neueren Meridianbeobachtungen die folgenden scheinbaren Positionen am Beobachtungstag zugrunde:

\begin{center}
\setlength{\tabcolsep}{3pt}
\begin{tabular}{llll}
Stern 1. & \shortstack[l]{Rektaszension\\$19^\circ\ 46'\ 4''0$} & \shortstack[l]{Deklination\\$0^\circ\ 56'\ 51''4$} & südlich \\
\hfill 2. & \hfill $19\ 52\ 18.8$ & \hfill $0\ 44\ 45.5$ & \\
\end{tabular}
\end{center}
\noindent Damit finde ich als Position des Kometen

\begin{center}
\setlength{\tabcolsep}{3pt}
\begin{tabular}{lll}
1844 Oktober 14. & $9^{\rm u}\ 27^{\rm m}\ 6^{\rm s}$ m. Z. & Rektasz. $20^\circ\ 16'\ 11''3$ \\
 & $9\ 32\ 7$ & Südl. Dekl. $0\ 44\ 30.3$ \\
\end{tabular}
\end{center}

\Needspace{7\baselineskip}
\SourceLeaf{475}
Dr. \textsc{Goldschmidt} hat für den Kometen folgende elliptische Elemente berechnet:

\begin{center}
Periheldurchgang 1844 Sept. 2.49152 mittlere Berliner Zeit.
\end{center}
{\small
\[
\left.\begin{array}{llr}
\text{Länge des Perihels} & \text{.\ .\ .\ .\ .} & 342^\circ\ 29'\ 44''9 \\
\text{Aufsteigender Knoten} & \text{.\ .\ .\ .\ .} & 63\ 48\ 55.2
\end{array}\right\}
\begin{array}{l}
\text{vom mittleren Äquinoktium} \\
\text{1844 Sept. 21.5 aus gezählt}
\end{array}
\]
}
\begin{center}
\setlength{\tabcolsep}{4pt}
\begin{tabular}{llr}
Bahnneigung & \mbox{.\ .\ .\ .\ .} & $2\ 55\ 1.9$ \\
\multicolumn{2}{l}{Logarithmus der großen Halbachse} & $0.4929151$ \\
Exzentrizität & \mbox{.\ .\ .\ .\ .\ .\ .\ .} & $0.6186103$ \\
\multicolumn{3}{l}{Rechtläufige Bewegung.}
\end{tabular}
\end{center}

\Needspace{14\baselineskip}
Die siderische Umlaufzeit, die sich aus diesen Elementen ergibt, beträgt 2400.339 Tage. Die Elemente geben die 3 zugrunde gelegten Beobachtungen (23. Aug., 21. Sept., 15. Okt.) genau wieder. Dr. \textsc{Goldschmidt} wird demnächst sämtliche bekannt gewordenen Beobachtungen mit diesen Elementen vergleichen. Vorläufig hat er die drei Beobachtungen vom 14. Okt., 1. Nov. und 2. Nov. damit verglichen. Dabei haben sich folgende Differenzen ergeben:

\begin{center}
\begin{tabular}{lrrl}
 & in Rektasz. & in Deklination & \\
Okt. 14 & $-0''9$ & $-8''3$ & Gauss \\
Nov. 1 & $+8.7$ & $-10.3$ & Petersen \\
\hfill 2 & $-12.3$ & $-8.4$ & Rümker \\
\end{tabular}
\end{center}

Die folgende Ephemeride ist nach diesen Elementen berechnet.

